\documentclass{scrartcl}
\usepackage[a4paper, total={7in, 10in}]{geometry}
\usepackage{stmaryrd}
\usepackage{amsthm}
\usepackage{amssymb}
\usepackage{amsmath}
\usepackage{algorithm2e}
\usepackage{hyperref}
\usepackage{comment}
\usepackage[french]{babel}
\usepackage{color}
\usepackage{tikz}
\usepackage{graphicx}
\usetikzlibrary{automata, arrows.meta, positioning}
\usepackage{MacrosArt}

%% Generated by ipescript page-labels figs.pdf
%% Do not edit
\newcommand{\ipeFigCliqueTree}{1}
\newcommand{\ipeFigCarteCliqueTree}{2}


\begin{document}

\title{Graphes de carte cordaux}

\author{José Lorgeré}

\maketitle

\begin{flushleft}

\section{Sous graphes interdits}

On note $K(3, G)$ le join d'un indépendant à $3$ sommets et du graphe $G$.

\begin{prop}\label{K3G}
    Soit $G = (A \sqcup B, E)$ un graphe tel que $A$ est composé de $3$ sommets
    indépendants de degré $3$, et tel que $B$ peut être contracté de telle sorte à ce que
    le graphe obtenu soit $K(3, G')$ pour un certain $G'$ à $3$ sommets. $G$ n'est pas de carte
\end{prop}

\begin{proof}
    On note $V = A \sqcup B$ et on se donne par l'absurde un témoin $H = (V \sqcup U, E')$
    de $G$. On va y construire un mineur $K_{3,3}$. Notons $v_1, v_2, v_3$ les sommets de $A$.
    Comme ils sont indépendants dans $G$ leur voisinages dans $H$ sont disjoints.
    \medskip \newline
    Maintenant considérons les $3$ ensembles de sommets de $B$ qui une fois contractés donnent
    transforment $G$ en $K(3, G')$. On les note $P_1, P_2, P_3$. Notons d'abord qu'un seul
    sommet par $P_i$ est relié à $v_j$, $i,j \in \{1, 2, 3\}$. Cela implique en particulier
    qu'une arrête interne à $P_i$ n'utilise pas de sommet de $N_H(v_j)$, pour être représentée
    dans $H$. Ainsi l'ensemble $N_H[P_j] \backslash N_H(A)$ est connexe, par connexité de
    $P_j$.
    \medskip \newline
    On contracte alors les ensembles suivants : $N_H[v_i]$, $N_H[P_j] \backslash N_H(A)$,
    $i,j \in \{1,2,3\}$. On note les sommets obtenus $V_i$ pour les premiers et $P'_j$ pour
    les seconds. Ces sommets forment alors un graphe $K_{3,3}$ absurde, comme $H$ était supposé
    planaire.
\end{proof}

Un cas particulier très utile de cette proposition est que le graphe $K(3, K_3)$ n'est pas
de carte. On notera $\mathcal{F}_3$ cette famille de graphes.
\newline \newline
On note $K^-_3(4, G)$ le graphe suivant : il s'agit du join d'un indépendant à $4$ sommets avec
$G$ à $4$ sommets, auquel on a retiré un couplage des indépendants vers $G$ de taille $3$.

\begin{prop}
    Soit $G = (A \sqcup B, E)$, où $A$ est composé de $4$ indépendants, 1 de degré $4$ et
    les autres de degré $3$, et tel que $B$ peut être contracté de manière à ce que le graphe
    obtenu soit un $K^-_3(4, G')$, pour $G'$ à $4$ sommets. $G$ n'est pas de carte.
\end{prop}

\begin{proof}
    Notons d'abord que le graphe $K^-_3(4, I_4)$ n'est pas planaire, où $I_4$ est un indépendant à
    $4$ sommets, car la planarité impliquerait que ce graphe a moins de $12$ arrêtes.
    \medskip \newline
    On suppose par l'absurde que $G$ est de carte, on s'en donne alors un témoin $H = (V \sqcup U, E')$,
    où $V = A \sqcup B$. On note $v_1, ..., v_4$ les $4$ sommets de $A$, où $v_4$ est le sommet
    de degré $4$. Comme pour les graphes précédents, les voisinages $N_H[v_i]$ sont disjoints.
    Notons $P_1, ..., P_4$ les ensembles de sommets contractés dans $B$ donnant un
    $K^-_3(4, G')$, où $P_4$ est le sommet ayant $4$ voisins dans $A$. La contrainte sur
    les degrés des sommets de $A$ donne là encore que $N_H[P_i] \backslash N_H(A)$ est connexe.
    On contracte alors ces ensembles de sommets et on obtient un mineur $K^-_3(4, I_4)$.
\end{proof}

On notera $\mathcal{F}_4$ cette famille de graphes.

\section{Approche par le graphe des cliques}

\subsection{Généralités sur les graphes cordaux}

On fixe $G$ un graphe cordal (connexe on va dire).

\begin{def*}
    Le graphe des cliques max de $G$ est le graphe $C(G) = (V_c, E_c)$ où $V_c$ est l'ensemble
    des cliques max de $G$, et où pour $C_1, C_2 \in V_c$, $C_1 C_2 \in E_c$ ssi $C_1 \cap C_2$
    est un séparateur minimal de $C_1 \backslash C_2$ et $C_2 \backslash C_1$
\end{def*}

On étiquette les arrête de ce graphe par le séparateur minimal obtenu par l'intersection de ses
deux extrémités.

\subsection{Contraintes sur les cartes cordales 3-connexes}

On suppose à présent que $G$ est cordal $3$-connexe et de carte.
On déduit de l'interdiction des sous graphes induit $\mathcal{F}_3$ et $\mathcal{F}_4$ les
contraintes suivantes sur le graphe des cliques max de $G$.

\begin{lem}\label{interlimit}
    Soient $C_1, C_2, C_3 \in V_c$ tels que $C_1C_2 \in E_c$. L'ensemble $C_3 \cap S_{12}$
    n'a pas plus de $2$ sommets.
\end{lem}

\begin{proof}
    Supposons par l'absurde que l'ensemble décrit ait au moins $3$ sommets. Montrons d'abord
    que $C_3 \not\subset (C_1 \cup C_2)$ : comme ce sont des cliques max, 
    $C_3 \backslash C_1$ et $C_3 \backslash C_2$ sont non vides. Si on avait l'inclusion,
    $S_{12}$ ne séparerait plus l'ensemble $C_1 \Delta C_2$, ce dernier étant connexe par
    la présence de $C_3$. \medskip \newline
    On choisit donc $v_3 \in C_3 \backslash (C_1 \cup C_2)$. On se donne alors
    $v_1 \in C_1$ tel que $v_1$ et $v_3$ soient indépendants, et $v_2 \in C_2$ vérifiant
    la même chose. $v_1$ et $v_2$ sont alors indépendants, car séparés par $S_{12}$. On
    forme alors un sous graphe induit $K(3, K_3)$ avec $S_{12}$ et ces $3$ indépendants, absurde.
\end{proof}

\begin{prop}
    Le graphe $C(G)$ est acyclique.
\end{prop}

\begin{proof}
    Supposons que l'on ait $C_1, ..., C_k, C_{k+1} = C_1$ un cycle de $C(G)$. Mais alors
    pour tout $1 \leq i \leq k$, $S_{i,i+1}$ sépare $C_i \Delta C_{i+1}$ dans $G$.
    Cela implique donc qu'il existe $j \neq i$ tel que $S_{j,j+1} \subset S_{i,i+1}$, le
    cas contraire on pourrait trouver un chemin en empruntant le cycle. On obtient alors
    une chaîne d'inclusions
    \[ S_{i_1, i_1+1} \supset S_{i_2, i_2+1} \supset ... \]
    et par principe des tiroirs on trouve $i \neq j$ tels que $S_{i,i+1} = S_{j,j+1}$.
    Les sommets $C_i, C_{i+1}, C_j$ (que l'on suppose sans perte de généralité distincts deux
    à deux) vérifient alors que
    \[ C_j \cap S_{i, i+1} = C_j \cap S_{j, j+1} = S_{j, j+1} \]
    Cet ensemble ayant plus de $3$ sommets comme $G$ est $3$-connexe, le Lemme \ref{interlimit}
    donne une contradiction.
\end{proof}

$C(G)$ est donc un arbre. Les résultats de (article clique graph) nous permettent alors
d'affirmer que pour tout sommet $v \in V$, l'ensemble des cliques max le contenant forme un
sous arbre de $C(G)$. On déduit de cela des résultats d'indépendances sur les cliques non
adjacentes dans l'arbre.

\begin{prop}\label{bornecliqueind}
    Soient $C_1, C_2 \in V_c$ tels que $C_1C_2 \notin E_c$. Alors $|C_1 \cap C_2| \leq 2$.
\end{prop}

\begin{proof}
    Considérons une arrête $C_i C_j \in E_c$ sur l'unique chemin de $C_1$ à $C_2$ dans $C(G)$.
    Les remarques précédentes sur les sous arbres de $C(G)$ permettent d'affirmer que
    $C_1 \cap C_2 \subset S_{ij}$.
    En appliquant le Lemme \ref{interlimit} à $C_1, C_i, C_j$ qu'on suppose distincts deux à
    deux quitte à échanger $C_1$ et $C_2$, on obtient ce qu'on veut.
\end{proof}

\begin{prop}\label{borneclique3}
    Soit $C^* \in V_c$ de degré au moins $3$ et considérons $C(G)$ comme enraciné en $C^*$,
    et notons $T_1, T_2, T_3$ trois sous arbres descendants de $C^*$. Pour tout triplet
    $C_i \in T_i$, $i \in \{1, 2, 3\}$, on a $|C_1 \cap C_2 \cap C_3| \leq 1$.
\end{prop}

% Reéxaminer avec attention plus tard

\begin{proof}
    Là encore la structure d'arbre des cliques contenant un sommet donné permet de se ramener
    au cas où $C_1, C_2, C_3$ sont racines de $T_1, T_2, T_3$, c'est à dire descendants directs
    de $C^*$.\newline
    Notons $S_i$ le séparateur $C_i \cap C^*$. $S_i$ sépare $C_i$ des deux autres $C_j$, $j \neq i$.
    En effet, $S_i \cap S_j$ a au plus $2$ sommets par le Lemme \ref{interlimit}, donc le retrait
    de $S_i$ ne sépare pas $C_j$ de $C^*$ car $G$ est $3$-connexe. Ainsi, si $C_i$ et $C_j$
    restaient connectés après le retrait de $S_i$, on pourrait construire un chemin de $C_i$ à $C^*$.
    \medskip \newline
    Par cette propriété, on choisit $v_i \in C_i$ non adjacent aux deux autres cliques pour
    tout $i \in \{1, 2, 3\}$. En supposans $|C_1 \cap C_2 \cap C_3| \geq 2$, on construit un
    sous graphe induit de $\mathcal{F}_3$ : les $v_i$ sont les indépendants, deux sommets
    leur étant complètement reliés sont les sommets de l'intersection, et on choisit un sommet
    par $S_i$ pour les $3$ derniers. C'est donc absurde et on a l'inégalité souhaitée.
\end{proof}

\begin{prop}\label{borneclique4}
    Soit $C^* \in V_c$ de degré au moins $4$, enracinons $C(G)$ en $C^*$ et considérons
    $T_1, T_2, T_3, T_4$ des sous arbres descendants de $C^*$. Pour tout quadruplet
    $C_i \in T_i$, si l'intersection de tout les $C_i$ est vide, alors pour tout
    $1 \leq i \leq 4$, il existe $k,l$ tels que $C_i \cap C_k \cap C_l = \varnothing$.
\end{prop}

\begin{proof}
    De manière similaire à la preuve précédente, on considère sans perte de généralité
    que $C_1, ..., C_4$ sont descendants directs de $C^*$. Supposons alors, sans perte
    de généralité, que l'intersection de ces $4$ cliques soit vides et que pour toute paire
    $k,l \neq 4$, on a $|C_4 \cap C_k \cap C_l| \geq 1$. Pour chaque telle paire $k,l$, on
    se donne un sommet $v_{k,l}$ dans cette intersection. Ils sont distincts deux à deux
    comme l'intersection totale est vide. Par des arguments similaires à la Proposition
    \ref{borneclique3}, on choisit également $v_i \in C_i$, $1 \leq i \leq 4$ les $4$
    étant indépendants.
    \medskip \newline
    On construit alors un sous graphe induit de $\mathcal{F}_4$ comme suit : les indépendants
    sont les $v_1, ..., v_4$. Les sommets de l'autre côté sont les $v_{k,l}$, chacun
    étant relié à $v_4, v_k, v_l$, auquels on ajoute un sommet par séparateur $S_i := C^* \cap C_i$,
    qui n'est alors adjacent qu'à $v_i$. Ces $4$ sommets peuvent être contractés en un seul,
    formant ainsi un $K^-_3(4, K_4)$. Absurde, on a ce qu'on veut.
\end{proof}

\subsection{Sommets longeants et graphe écorce}

On suppose ici uniquement que $G$ est cordal $3$-connexe et que $C(G)$ est un arbre.

\begin{def*}
    Un sommet $v \in V$ tel qu'il existe $C_1, C_2 \in V_c$, $C_1C_2 \notin E_c$ tels que
    $v \in C_1 \cap C_2$ sera dit longeant.
\end{def*}

Là encore par les résultats sur les clique tree, l'ensemble des sommets de $C(G)$ contenant
un sommet longeant forme un arbre, qui dans ce cas possède au moins $3$ sommets. Pour un sommet
longeant $v$, on notera $T_v$ l'arbre des cliques le contenant. Les
résultats obtenus dans la section précédente donnent donc des bornes sur le nombre de sommets
longeants contenus dans les cliques de $G$.

\begin{def*}
    On définit le graphe écorce de $G$, noté $B(G)$ comme le graphe obtenu depuis $C(G)$ comme suit :
    pour tout sommet longeant $v \in V$, on l'ajoute au graphe et on le relie à $T_v$.
\end{def*}

On dispose du lemme technique suivant concernant les sommets longeants de $B(G)$.

\begin{lem}\label{voislongpath}
    Soit $v \in V$ un sommet longeant, relié à $C_1, C_2, C_3 \in V_c$ dans $B(G)$. Supposons
    qu'il existe pour tout $i \in \{1,2,3\}$ un chemin de $C_i$ à $C'_i \in V_c$ tels que les
    deux chemins soient disjoints. Alors il existe $C''_1, C''_2, C''_3 \in V_c$ voisins de $v$
    formant un chemin et vérifiant la même propriété relativement aux $C'_i$.
\end{lem}

\begin{proof}
    Notons $P_i$, $i \in \{1, 2, 3\}$ les $3$ chemins disjoints reliant $C_i$ à $C'_i$.
    $C(G)$ étant un arbre, il existe un unique chemin $P_{ij}$ reliant $C_i$ à $C_j$,
    $i \neq j$. $P_{ij}$ rencontre alors $P_j$. On suppose sans perte de généralité
    comme $C(G)$ est acyclique que
    $P_{12}$ ne rencontre pas $P_3 \backslash P_2$ et que $P_{32}$ ne rencontre pas $P_1 \backslash P_2$.
    \medskip \newline
    Notons $C''_2$ la dernière intersection (dans l'ordre de $C_2$ à $C'_2$) entre $P_2$
    et un des deux autres chemins. On suppose sans perte de généralité qu'il s'agit de
    celle avec $P_{12}$. On définit alors $C''_1$, le prédecesseur de $C''_2$ dans
    $P_{12}$ et $C''_3$ le prédecesseur de $C''_2$ dans $P_{32}$. Les trois forment alors
    bien un chemin. On peut également construire des chemins disjoints menant chacun d'eux
    aux $C'_i$ correspondant. Enfin $T_v$ est un arbre contenant $C_1, C_2, C_3$, donc
    contient $C''_1, C''_2, C''_3$.
\end{proof}

% Reformuler

On cherche alors à relier la planarité de $B(G)$ au fait que $G$ soit de carte. On introduit
pour cela la notion d'ensembles non planaires minimaux : notons $L \subset V$ l'ensemble
des sommets longeants. Un sous ensemble $S \subset L$ sera dit non planaire minimal si
$B(G) - (L\backslash S)$
n'est pas planaire et que $S$ est minimal pour cette propriété. On utilisera pour toute
la suite la caractérisation de cette propriété par mineurs interdits.
\medskip \newline
On rappelle qu'un graphe $H = (\mathcal{V}, \mathcal{E})$ est mineur de $B(G)$ si et seulement
si à isomorphisme près, on a $\mathcal{V} = \{T_1, ..., T_r\}$, des sous arbres disjoints
de $B(G)$ tels que $T_i T_j \in \mathcal{E}$ si et seulement si il existe un sommet
de $T_i$ et un de $T_j$ reliés par une arrête dans $B(G)$. Il se trouve que l'on peut requérir
dans le cas qui nous intéresse un choix particulier d'arbres.

\begin{lem}\label{normalformK33}
    Soit $S$ un ensemble non planaire minimal donnant lieu à un mineur $K_{3,3}$. Ce dernier
    peut être choisi de telle sorte à ce que tout les arbres d'un côté de la bipartition
    contiennent au plus un sommet longeant.
\end{lem}

\begin{proof}
    On montre seulement que l'on peut toujours déplacer un sommet longeant d'un arbre à un
    autre adjacent, si le premier a en a au moins $2$. On se donne un mineur $K_{3,3}$ ou
    intervient $S$. On écrit $T_1, T_2, T_3$ et $T'_1, T'_2, T'_3$ les arbres des deux côtés
    de la bipartition. Supposons sans perte de généralité que $T_1$ possède un sommet longeant.
    Soient $p_1 \in T_1$, $q_1 \in T'_1$ tels que $p_1q_1$ soit une arrête de $B(G)$. On se
    donne $v_1 \in T_1$ longeant tel que le chemin de $p_1$ à $v_1$ soit minimal pour
    l'inclusion. On le note $P_1$. De même, on se donne $p_2$ et $v_2$ relatifs à $T'_2$ et
    on note $P_2$ le chemin minimal.
    \medskip \newline
    Montrons que l'on peut soit ajouter $P_1$ à $T'_1$ et le retirer à $T_1$, soit
    faire de même pour $P_2$ et $T'_2$ tout en conservant la structure du graphe. Supposons
    par l'absurde qu'on ne puisse pas. Alors tout les sommets adjacents à $T'_3$ dans $T_1$
    sont dans $P_1 \cap P_2$. Si l'on remplace $T_1$ par $P_1 \cup P_2$, on obtient un arbre
    contenant un unique sommet longeant et adjacent à $T'_1, T'_2, T'_3$. Absurde,
    on a construit un mineur $K_{3,3}$ sans utiliser un sommet de $S$, ce qui contredit la
    minimalité.
\end{proof}

% \begin{lem}
%     Soit $S$ un ensemble non planaire minimal donnant lieu à un mineur $K_5$. Ce dernier peut
%     être choisi de telle sorte à ce que si un arbre possède plus de deux sommets longeants,
%     au plus un arbre n'en possède aucun.
% \end{lem}

\begin{lem}\label{minimalmax3}
    Un ensemble non planaire minimal a au plus $3$ sommets.
\end{lem}

\begin{proof}
    Trouver une preuve non abominable (ou tout du moins raisonnable).
\end{proof}

\begin{prop}
    Si $G$ est de carte, alors $B(G)$ est planaire.
\end{prop}

\begin{proof}
    On raisonne par contraposée. On suppose donc $B(G)$ non planaire, c'est à dire contenant
    un mineur $K_{3,3}$ ou $K_5$, et on veut en déduire des structures interdites
    pour un graphe de carte.
    \newline \newline
    Si $B(G)$ est non planaire, l'ensemble $L$ des sommets longeants est un ensemble non planaire.
    On se donne alors $S$ un sous ensemble non planaire minimal. On distingue selon le mineur
    produit. On s'appuie également sur le Lemme \ref{minimalmax3} pour simplifier l'étude,
    et le fait qu'un ensemble non planaire minimal a au moins $2$ sommets.
    \begin{itemize}
        \item Cas $K_{3,3}$ : on écrit $T_1, T_2, T_3$ et $T'_1, T'_2, T'_3$ les arbres des
        deux côtés de la bipartition, donnant lieu au mineur.
        \begin{itemize}
            \item[Si $|S| = 2$ :] Notons que par acyclicité de $C(G)$, les deux sommets
            longeants sont du même côté de la bipartition dans des arbres différents. Dans
            les deux cas, par le Lemme \ref{normalformK33} on peut construire un mineur
            $K_{2,2}$ n'utilisant aucun sommet longeant, qui donne alors lieu à un cycle dans
            $C(G)$, absurde. Pour le cas restant, notons $v_1, v_2$ les deux sommet longeants
            respectivement dans $T_1$ et $T_2$, wilog. On choisit $C^* \in T_3$ de degré
            au moins $3$ et
            $C_1, C_2, C_3$ respectivement dans $T'_1, T'_2, T'_3$, arbres que l'on
            enracine en ces sommets. On considère le chemin $P^{(j)}_i$ dans
            $T_j \cup T_i$ de $v_j$ à $C_i$, $j \in \{1,2\}, i \in \{1,2,3\}$. $v_1$ a alors
            pour voisin les extrémités des $P^{(j)}_i$, et en exploitant les chemins
            de $C^*$ vers les $C_i$ et le fait que $v_j$ est adjacent à un sous arbre,
            on a une arrête directe entre $v_j$ et chaque $C_i$. Par définition de $B(G)$,
            $|C_1 \cap C_2 \cap C_3| \geq 2$. La Proposition \ref{borneclique3} donne
            alors que $G$ n'est pas de carte.

            \item[Si $|S| = 3$ :] Par minimalité de $S$, les trois sommets longeants sont
            répartis un par $T_i$. Sinon on obtient wilog par le Lemme \ref{normalformK33}
            $v_1, v_2$ dans $T_1, T_2$ longeants et $v_3$ dans $T'_3$. Les chemins de $v_i$
            à $T'_j$, $i,j\in \{1,2\}$ permettent d'affirmer par la structure de $B(G)$
            qu'il existe $C^{(j)}_1, C^{(j)}_2$ respectivement dans $T'_1, T'_2$ voisins directs
            de $v_j$. Notons $C'_1, C'_2, C'_3$ les voisins de $v_3$ sur le chemin de $T_1, T_2, T_3$
            respectivement. En exploitant la structure d'arbre des voisins de $v_1, v_2$, et les chemins
            entre $C \in T_3$ et $T'_1, T'_2, T'_3$, on obtient que $v_1, v_2$ sont adjacents à $C'_1, C'_2, C'_3$,
            et on peut ainsi, comme il existe des chemins entre ces derniers, retirer $v_3$ et garder
            un mineur $K_{3,3}$.
            \medskip \newline
            On a donc bien la situation décrite au départ. On exploite ENCORE la structure d'arbre
            des voisins des sommets longeants pour montrer qu'on a en fait un sous graphe $K_{3,3}$
            et non pas un mineur (les arguments sont vraiment presque les mêmes). Parmi
            les $3$ sommets de ce sous graphe, au moins $2$ ne sont pas adjacents, par acyclicité
            de $C(G)$, mais partagent $3$ sommets longeants. La Proposition \ref{bornecliqueind}
            donne que $G$ n'est pas de carte.
        \end{itemize}

        \item Cas $K_5$ : on écrit $T_1, ... T_5$ les sous arbres donnant lieu à ce mineur.
        Dans ce cas, il est clair que $|S| = 3$ et que les sommets longeants sont répartis un
        par arbre,car on aurait un cycle dans $C(G)$ le cas contraire. On suppose $v_1, v_2, v_3$
        dans $T_1, T_2, T_3$. En exploitant les mêmes arguments que précédemment, on peut montrer
        qu'il existe $C_4, C_5 \in T_4, T_5$ respectivements tels que $v_1, v_2, v_3$ soient
        voisins directs de ces derniers.\newline
        XXXXX\newline
        Pour la suite, on va exploiter les chemins entre $v_i$ pour exhiber une structure
        $\mathcal{F}_4$ donnée par une des propositions précédentes.
    \end{itemize}
\end{proof}

\subsection{Construction de cartes 3-connexes}

On se fixe à présent $G$ un graphe cordal $3$-connexe, dont on suppose que $C(G)$ est un arbre.

\begin{prop}\label{carte3conndisj}
    Si $G$ n'a pas de sommets longeants, alors $G$ est de carte.
\end{prop}

\begin{proof}
    On construit un témoin de $G$ comme suit : on pose $H = (V \sqcup V_c, E')$, où pour tout
    $v \in V$ on relie $v$ à tout les $C \in V_c$ le contenant. Ce graphe est planaire :
    on peut partir d'un plongement de $C(G)$ dans le plan pour en construire un de $H$, il
    suffit pour cela de placer chaque $V_c$ a la position prescrite par le plongement de $C(G)$
    et de disposer les sommets de $V$ comme illustré Figure \ref{figcarte3conndisj} : les sommets faisant partie
    d'une seule clique max sont placés transversalement aux arrêtes de $C(G)$, tandis que
    ceux étant reliés à $2$ cliques max sont placés parallélement à ces dernières. Un sommet
    est au plus relié à $2$ cliques max comme $G$ n'a pas de sommets longeants.
\end{proof}

\begin{figure}[h]
    \caption{Illustration de la Proposition \ref{carte3conndisj}}\label{figcarte3conndisj}
    \begin{center}
    A gauche un plongement de $C(G)$, à droite le plongement du témoin $H$.\\~\\
    \includegraphics*[page=\ipeFigCliqueTree, scale = 0.60]{figscordal}
    \hspace*{0.5cm}
    \includegraphics*[page=\ipeFigCarteCliqueTree, scale = 0.60]{figscordal}
    \end{center}
\end{figure}

Reste à traiter le cas où des cliques non reliées par une arrête dans $C(G)$ se rencontrent.
\medskip \newline
On construit alors le graphe $C'(G)$ depuis $C(G)$ comme suit : pour tout sommet longeant $v$,
on ajoute ce dernier au graphe et on le relie à tout les sommets de $T_v$. On a alors la
caractérisation suivante pour $G$.
\begin{prop}
    $G$ admet une carte si et seulement si $C'(G)$ est planaire.
\end{prop}

\begin{proof}
    On a déjà montré un sens. Supposons que $C'(G)$ soit planaire. On se donne un plongement
    de ce dernier. On construit d'abord un témoin de $G - L$, où $L$ est l'ensemble des
    sommets longeants, par la procédure donnée à la Proposition \ref{carte3conndisj}. Ensuite,
    on ajoute les sommets longeants : on ajoute au témoin les sommets longeants de $G$,
    et on ajoute des arrêtes entre eux et chaque clique maximale à laquelle ils appartiennent.
    Le plongement de ces arrêtes du témoin est donné par le plongement planaire
    de $B'(G)$.
\end{proof}

\end{flushleft}

\end{document}